% Définition diagrammatique de la matrice T : développement complet (première
% ligne) et forme récursive de Dyson (deuxième ligne).
% Style identique aux diagrammes du chapitre Raman (tikz/e1.tex, tikz/pd.tex).
% Compatible pdflatex. Requiert : tikz, decorations.markings, arrows.meta, calc.
% Le blob hachuré est tracé à la main (clip + traits) pour éviter d'avoir à
% charger la bibliothèque patterns dans le préambule.

% Vertex électron-défaut : cercle blanc + croix (cf. tikz/règles_Raman.tex)
\providecommand{\vertexdefaut}[1]{%
  \draw[fill=white] (#1) circle (0.2);
  \draw[line width=1pt] ($(#1)+(-0.14,-0.14)$) -- ($(#1)+(0.14,0.14)$);
  \draw[line width=1pt] ($(#1)+(-0.14,0.14)$) -- ($(#1)+(0.14,-0.14)$);
}

% Matrice T : cercle hachuré. #1 = coordonnée du centre.
\providecommand{\blobT}[1]{%
  \begin{scope}
    \clip (#1) circle (0.28);
    \foreach \i in {-8,...,8} {
      \draw[line width=0.4pt] ($(#1)+(0.08*\i-0.45,-0.45)$) -- ($(#1)+(0.08*\i+0.45,0.45)$);
    }
  \end{scope}
  \draw[thick, fill=none] (#1) circle (0.28);
}

\begin{tikzpicture}[
    line/.style={thick},
    >=Stealth,
    fermion/.style={
      postaction={decorate},
      decoration={markings, mark=at position 0.55 with {\arrow{Stealth}}}
    },
]

% ============ première ligne : développement complet ============
\node[anchor=east] at (1.5,0) {$T_{mn} \equiv$};
\blobT{1.9,0}
\node at (2.6,0) {$=$};

% un vertex
\vertexdefaut{3.3,0}
\node at (3.9,0) {$+$};

% deux vertex
\vertexdefaut{4.5,0}
\draw[line, fermion] (4.7,0) -- (5.7,0);
\vertexdefaut{5.9,0}
\node at (6.5,0) {$+$};

% trois vertex
\vertexdefaut{7.1,0}
\draw[line, fermion] (7.3,0) -- (8.3,0);
\vertexdefaut{8.5,0}
\draw[line, fermion] (8.7,0) -- (9.7,0);
\vertexdefaut{9.9,0}
\node[anchor=west] at (10.35,0) {$+\ \cdots$};

% ============ deuxième ligne : forme récursive ============
\node at (2.6,-1.7) {$=$};

\vertexdefaut{3.3,-1.7}
\node at (3.9,-1.7) {$+$};

\vertexdefaut{4.5,-1.7}
\draw[line, fermion] (4.7,-1.7) -- (5.72,-1.7);
\blobT{6.0,-1.7}

\end{tikzpicture}
